\section{Measuring Deployed Agent Runtimes}\label{sec:audit}

The attacks of \cref{sec:attacks} ask whether a design can be broken. The audit asks a narrower question of
the tools developers use: when a runtime pauses for human approval, which of \Prop{1} to \Prop{5} and
\Prop{7} does that pause have?

\subsection{Method}

We wrote one probe per runtime, ran each in its own isolated environment with a stub model so that results do not
depend on a network model, and recorded every outcome together with an evidence log. We executed probes against nine
runtimes: the OpenAI Agents SDK, LangGraph, Google ADK, Agno, Pydantic AI, CrewAI, AutoGen AgentChat, smolagents
and the MCP Python SDK. The tenth, OpenClaw, needs a model for its live loop, so we reviewed it statically from
its released bundle. The smolagents library has no built-in approval mechanism, so every cell of its row is n/a. For CrewAI and AutoGen we audited both the general human-input primitive and the documented
approval gate, which gives twelve rows. Claude Code, Codex CLI and Cursor are proprietary; we wrote manual test
procedures for them but did not run them, and the table excludes them. Three rules keep the criteria uniform
across runtimes.

\paragraph{\Prop{1} is reported as structural, not measured.} An approval callback that runs in the agent's own
process is reachable by the agent by construction under this threat model. We mark those cells $\dagger$ and
exclude them from every count, and instead measure two observable facts: whether a paused run can be resumed from
another process, and whether the runtime ships an approver that is a separate principal.

\paragraph{\Prop{2} uses a capability ladder.} Level L1 is the approver editing arguments, which is a legitimate
channel and not a gap. Level L2 is any party that can call the documented resume or continue API. Level L3 is a
writer of the backing store, such as a checkpoint or session database. A runtime passes if it resists L2 and L3,
is partial if it resists only L2, and fails if an L2 party can change what executes after approval.

\paragraph{\Prop{3} uses one replay criterion.} For every runtime we tried a same-call replay, a rollback of run
state followed by a replay, and a cross-call replay. A runtime whose approvals are consumed relative to run state
but double-execute after a rollback is marked partial.

Cells whose outcome is a reading of source or documentation rather than an observed value are marked $J$
(judgment); all \Prop{4} and \Prop{5} cells are of this kind, because no experiment can select an outcome without
a reference renderer or verifier in each runtime. There are \auJudgement{} judgment cells, and they are excluded
from the counts below.

\begin{table}[t]
\centering\small
\caption{Approval properties of agent runtimes. Subscripts: \textsubscript{state}, single use holds only relative
to run state, and a rollback replays; \textsubscript{store}, binding resists the resume API but not a writer of the
backing store; \textsubscript{int}, holds only if the integrator renders from structured data;
\textsubscript{low}, low-severity gap. Superscripts: $\dagger$, fails by construction under the threat model;
$J$, judgment from source or API review; $s$, static review only. n/a: no such mechanism. \Prop{7} is
measured within one run.}
\label{tab:audit}
\setlength{\tabcolsep}{4pt}
\begin{tabular}{@{}l l c c c c c c@{}}
\toprule
Runtime & Version & \Prop{1} & \Prop{2} & \Prop{3} & \Prop{4} & \Prop{5} & \Prop{7}\\
\midrule
\inputtab{tables/tab_audit.tex}
\bottomrule
\end{tabular}
\end{table}

\subsection{Results}

\Cref{tab:audit} gives the outcome for every runtime. We highlight four findings.

\paragraph{Effect binding splits the field.} Of the eight rows where binding could be measured, \auPtwoPass{} pass,
\auPtwoPart{} is partial and \auPtwoFail{} fail. The OpenAI Agents SDK binds a call identifier to one canonical
payload and rejects a resumed run whose arguments differ. AutoGen's code-executor approval executes exactly the
code that was approved, and MCP keeps arguments on the server. Google ADK resisted all five mismatched
confirmations we sent through its runner, but a writer of the session store can substitute arguments. LangGraph and
Agno execute whatever arguments are in thread state at resume time, so a party holding the thread substitutes the
effect with a bare ``approve''. Pydantic AI executes the arguments in a caller-supplied message history, and a later
CrewAI hook can mutate arguments that an earlier hook approved.

\paragraph{The substitution persists past its published range.} Loopjacking reports post-approval substitution in
Agno's AgentOS for versions 2.5.6 to 3.0.9~\citep{arunkumar2026loopjacking}. We re-ran the same probe over real HTTP
against 3.0.9, 3.0.10 and the latest release, 3.1.3 (\cref{tab:agno}). Substitution through the documented remote
continue endpoint reproduces on all three. Same-call replay, which double-executed on the 3.0 releases, executes
once on 3.1.3, while a rollback still double-executes. The attack class is already public. Only the affected version range is new, and we make no novelty claim for the
class.

\begin{table}[t]
\centering\small
\caption{Version dependence in Agno. The same probe, with the same criteria, re-run in each release. 3.1.3 is the
latest release measured.}
\label{tab:agno}
\begin{tabular}{@{}l c l l l l@{}}
\toprule
Agno & In published range & \Prop{2} & Remote continue & \Prop{3} & Same-call replay\\
\midrule
3.0.9 & yes & FAIL & substituted & FAIL & re-executes\\
3.0.10 & no & FAIL & substituted & FAIL & re-executes\\
3.1.3 & no & FAIL & substituted & PART\textsubscript{state} & consumed once\\
\bottomrule
\end{tabular}
\end{table}

\paragraph{Single use is relative to run state in most runtimes.} In \auPthreePart{} rows, an approval is consumed
relative to the run but executes twice when the run state is rolled back and resumed, because the approval artifact carries no
nonce or expiry of its own. The two rows that pass, the CrewAI and AutoGen gates, ask again for every call.

\paragraph{No measured runtime limits floods, and none produces a verifiable record.} In all \auPsevenFail{} rows where flooding
could be measured, nothing capped or rate-limited pending approvals. In the five runtimes with a native
pause, a single model turn raised 20 pending approvals. By our reading of their APIs, no runtime produces a
signed approval record that a third party could verify. Approvals are stored as unsigned booleans or flags in
run state. Every runtime with structured arguments leaves rendering to the integrator, and only OpenClaw ships a card
that escapes invisible characters and caps length.

These results describe each runtime's documented approval primitive, not every way an integrator might use it. A
team that renders arguments carefully, persists approvals with a nonce, and resumes runs only from a trusted
service can close several of these gaps. The point of the audit is that, by default, the pause a developer adds
is not yet a factor the agent cannot reach.
